\documentclass{beamer}

% Configuración del tema
\usetheme{Madrid}
\usecolortheme{default}

% Paquetes para español y codificación
\usepackage[utf8]{inputenc}
\usepackage[T1]{fontenc}
\usepackage[spanish]{babel}
\usepackage{graphicx}
\usepackage{tikz}
\usetikzlibrary{positioning, arrows.meta, shapes.geometric}
\usepackage{hyperref}

% Información del título
\title[Clase 1: Fundamentos y Formulación]{Clase 1: Fundamentos de la Web y Formulación del Proyecto}
\subtitle{Proyecto Web}
\author{Rodrigo García, PhD}
\date{\today}

\begin{document}

% Diapositiva de Título
\begin{frame}
    \titlepage
\end{frame}

% Índice
\begin{frame}{Agenda de la Clase}
    \tableofcontents
\end{frame}

% ==========================================
% SECCIÓN 1: CONCEPTOS BÁSICOS DE LA WEB
% ==========================================
\section{Conceptos Básicos de la Web}

\begin{frame}{¿Qué es la Web?}
    \begin{itemize}
        \item La World Wide Web (WWW) es un sistema de información interconectado accesible a través de Internet.
        \item Diferencia clave: Internet es la infraestructura (cables, routers), la Web es el servicio que usamos sobre ella.
        \item Creada por Tim Berners-Lee en 1989.
    \end{itemize}
\end{frame}

\begin{frame}{Arquitectura Cliente-Servidor (Visualización)}
    \begin{center}
    \begin{tikzpicture}[
        node distance=4.5cm,
        box/.style={draw, rounded corners, align=center, inner sep=2mm, fill=blue!10, minimum height=1.5cm, minimum width=2.5cm},
        arrow/.style={-{Stealth[length=2mm]}, thick},
        textlabel/.style={font=\small, align=center}
    ]
        
        % Nodos
        \node[box] (cliente) {\textbf{Cliente}\\(Navegador Web)};
        \node[box, right=of cliente, fill=green!10] (servidor) {\textbf{Servidor}\\(Aplicación / BD)};
        
        % Flechas
        \draw[arrow, color=blue!70!black] ([yshift=3mm]cliente.east) -- ([yshift=3mm]servidor.west) 
            node[midway, above, textlabel] {1. \textbf{Petición HTTP/HTTPS}\\(Ej: GET /index.html)};
            
        \draw[arrow, color=green!60!black] ([yshift=-3mm]servidor.west) -- ([yshift=-3mm]cliente.east) 
            node[midway, below, textlabel] {2. \textbf{Respuesta HTTP/HTTPS}\\(Ej: Envía archivo HTML)};
            
    \end{tikzpicture}
    \end{center}
    
    \vspace{0.3cm}
    \begin{itemize}
        \item \textbf{Cliente:} Inicia la comunicación solicitando información.
        \item \textbf{Servidor:} Escucha peticiones, procesa la lógica y devuelve resultados.
        \item \textbf{Protocolo HTTP/HTTPS:} Es el "idioma" estándar de la web. HTTPS incluye una capa de cifrado para seguridad.
    \end{itemize}
\end{frame}

\begin{frame}{La Tríada de la Web Frontend (Visualización)}
    \begin{center}
    \begin{tikzpicture}[
        node distance=0.1cm and 0.1cm,
        box/.style={draw, rounded corners, align=center, inner sep=2mm, minimum height=2.2cm, minimum width=3cm},
        html/.style={box, fill=orange!20, draw=orange!80!black, thick},
        css/.style={box, fill=blue!20, draw=blue!80!black, thick},
        js/.style={box, fill=yellow!30, draw=yellow!80!black, thick}
    ]
        
        \node[html] (html) {\Large\textbf{HTML}\\[1ex] \textbf{Estructura}\\[0.5ex] \small(El esqueleto)};
        \node[right=0.1cm of html] (plus1) {\Large \textbf{+}};
        
        \node[css, right=0.1cm of plus1] (css) {\Large\textbf{CSS}\\[1ex] \textbf{Estilo}\\[0.5ex] \small(La piel/ropa)};
        \node[right=0.1cm of css] (plus2) {\Large \textbf{+}};
        
        \node[js, right=0.1cm of plus2] (js) {\Large\textbf{JS}\\[1ex] \textbf{Acción}\\[0.5ex] \small(Los músculos)};
            
    \end{tikzpicture}
    \end{center}
    
    \vspace{0.3cm}
    \begin{itemize}
        \item \textbf{HTML:} Define qué elementos existen en la página (párrafos, imágenes, botones, formularios).
        \item \textbf{CSS:} Define cómo se ven visualmente esos elementos (colores, posición, tipografía, responsividad).
        \item \textbf{JavaScript:} Define qué ocurre al interactuar (animaciones, clics, carga de datos dinámicos).
    \end{itemize}
\end{frame}

\begin{frame}[fragile]{Profundizando en HTML: La Estructura}
    \begin{itemize}
        \item \textbf{Significado:} \textit{HyperText Markup Language} (Lenguaje de Marcado de Hipertexto).
        \item \textbf{Propósito:} Organizar y estructurar el contenido de la web usando etiquetas (tags).
    \end{itemize}
    
    \vspace{0.3cm}
    \textbf{Ejemplo de sintaxis:}
    \begin{verbatim}
    <h1>Bienvenido a mi Proyecto</h1>
    <p>Este es un texto descriptivo del problema.</p>
    <a href="https://ejemplo.com">Enlace de prueba</a>
    <img src="logo.png" alt="Logotipo">
    \end{verbatim}
    
    \vspace{0.3cm}
    \begin{itemize}
        \item Las etiquetas suelen tener apertura \verb|<etiqueta>| y cierre \verb|</etiqueta>|.
        \item \textbf{Regla de oro:} HTML debe centrarse en \textbf{QUÉ} es el contenido (semántica), no en \textbf{CÓMO} se ve visualmente.
    \end{itemize}
\end{frame}

\begin{frame}{Backend, BD y API (Visualización)}
    \begin{center}
    \begin{tikzpicture}[
        node distance=1.5cm and 1.8cm,
        box/.style={draw, rounded corners, align=center, inner sep=2mm, minimum height=1.5cm, minimum width=2.2cm},
        front/.style={box, fill=orange!20, draw=orange!80!black, thick},
        api/.style={box, fill=gray!20, draw=gray!80!black, thick, dashed},
        back/.style={box, fill=green!20, draw=green!80!black, thick},
        db/.style={box, cylinder, shape border rotate=90, aspect=0.25, fill=purple!20, draw=purple!80!black, thick, minimum height=1.8cm, minimum width=1.6cm},
        arrow/.style={-{Stealth[length=2mm]}, thick},
        bidir/.style={<->, >=Stealth, thick},
        textlabel/.style={font=\scriptsize, align=center}
    ]
        
        % Nodos
        \node[front] (frontend) {\textbf{Frontend}\\(Navegador)};
        \node[api, right=of frontend] (api) {\textbf{API}\\(Puente)};
        \node[back, right=of api] (backend) {\textbf{Backend}\\(Lógica)};
        \node[db, below=of backend, yshift=0.5cm] (database) {\textbf{Base de}\\\textbf{Datos}};
        
        % Flechas
        \draw[bidir, color=blue!70!black] (frontend.east) -- (api.west) 
            node[midway, above, textlabel] {Petición JSON}
            node[midway, below, textlabel] {Respuesta};
            
        \draw[arrow, color=gray!80!black] (api.east) -- (backend.west);
        
        \draw[bidir, color=red!70!black] (backend.south) -- (database.north) 
            node[midway, right, textlabel] {SQL / NoSQL};
            
    \end{tikzpicture}
    \end{center}
    
    \vspace{0.2cm}
    \begin{itemize}
        \item \textbf{Backend:} Contiene la lógica segura. Recibe órdenes, procesa datos y devuelve resultados (Node.js, Python, Java).
        \item \textbf{BD (Base de Datos):} Almacén estructurado y persistente (PostgreSQL, MongoDB).
        \item \textbf{API:} El "menú". Define las reglas y formatos en los que el Frontend puede pedir datos al Backend.
    \end{itemize}
\end{frame}

% ==========================================
% SECCIÓN 2: IDENTIFICACIÓN DEL PROBLEMA
% ==========================================
\section{Identificación del Problema y Usuarios Objetivo}

\begin{frame}{El Proyecto Integrador}
    \begin{block}{¿Qué vamos a construir?}
        Una aplicación web completa a lo largo del semestre que integre Frontend, Backend, Base de Datos, despliegue en la Nube y medidas de Seguridad.
    \end{block}
    \vspace{0.5cm}
    Para que el proyecto tenga sentido, debe resolver un \textbf{problema real}.
\end{frame}

\begin{frame}{Identificación del Problema}
    \begin{itemize}
        \item Todo buen proyecto web nace de una necesidad.
        \item \textbf{Preguntas clave:}
        \begin{itemize}
            \item ¿Qué proceso es actualmente ineficiente o manual?
            \item ¿Qué información es difícil de encontrar o gestionar?
            \item ¿Quién está sufriendo este dolor (pain point)?
        \end{itemize}
        \item \textbf{Ejemplos:} Gestión de inventario para un negocio local, plataforma de reservas para servicios, sistema de control para grupos estudiantiles.
    \end{itemize}
\end{frame}

\begin{frame}{Definición de Usuarios Objetivo}
    \begin{itemize}
        \item ¿Para quién estamos construyendo esto?
        \item No todos los usuarios son iguales (Roles: Administrador, Cliente, Empleado).
        \item \textbf{User Personas:} Crear perfiles ficticios pero realistas de quién usará la app.
        \item Entender el contexto técnico del usuario (¿Usan móvil o PC? ¿Tienen buena conexión a internet?).
    \end{itemize}
\end{frame}

% ==========================================
% SECCIÓN 3: EQUIPOS Y REPOSITORIO
% ==========================================
\section{Equipos de Trabajo y Control de Versiones}

\begin{frame}{Formación de Equipos}
    \begin{itemize}
        \item El trabajo en el desarrollo web es altamente colaborativo.
        \item \textbf{Tamaño ideal:} 4 a 5 personas.
        \item \textbf{Roles sugeridos (aunque todos toquen código):}
        \begin{itemize}
            \item Frontend Lead
            \item Backend/API Lead
            \item Scrum Master / Project Manager (Organización)
            \item DevOps / Cloud Lead (Despliegue)
        \end{itemize}
        \item \textbf{Dinámica:} Formación de grupos y selección de una idea de proyecto.
    \end{itemize}
\end{frame}

\begin{frame}{Responsabilidades por Rol}
    \begin{itemize}
        \item \textbf{Frontend Lead:} Define la arquitectura de la interfaz. Vela por la experiencia de usuario (UX), accesibilidad y el diseño responsivo.
        \item \textbf{Backend/API Lead:} Diseña la base de datos, crea los endpoints de la API y asegura que la lógica de negocio y la seguridad estén correctamente implementadas.
        \item \textbf{Scrum Master / Project Manager:} Facilita el trabajo del equipo, organiza las tareas en el tablero Kanban (GitHub Projects, Trello), y se asegura de cumplir los hitos de entrega.
        \item \textbf{DevOps / Cloud Lead:} Configura y mantiene el repositorio (Git/GitHub), gestiona el flujo de trabajo (ramas, PRs), y despliega la aplicación en servicios en la nube.
    \end{itemize}
\end{frame}

\begin{frame}{Control de Versiones (Git y GitHub)}
    \begin{itemize}
        \item ¿Qué pasa si dos personas editan el mismo archivo al mismo tiempo? Git lo resuelve.
        \item \textbf{Git:} Sistema de control de versiones local.
        \item \textbf{GitHub/GitLab:} Plataforma en la nube para alojar repositorios Git y colaborar.
        \item \textbf{Conceptos básicos:} \texttt{commit}, \texttt{push}, \texttt{pull}, \texttt{branch}, \texttt{merge}.
    \end{itemize}
\end{frame}

\begin{frame}{Creación del Repositorio del Proyecto}
    \begin{block}{Actividad Práctica}
        \begin{enumerate}
            \item Crear una organización o repositorio en GitHub.
            \item Invitar a todos los miembros del equipo.
            \item Configurar el archivo \texttt{README.md} inicial con el nombre del proyecto y la descripción del problema.
            \item (Opcional) Configurar el archivo \texttt{.gitignore}.
        \end{enumerate}
    \end{block}
\end{frame}

\begin{frame}{Ejemplos de Proyectos: Problema y Usuarios}
    \begin{small}
    \begin{itemize}
        \item \textbf{1. Periódico Escolar Digital:} \\ \textit{Problema:} Los estudiantes no se enteran de las noticias o eventos del colegio. \\ \textit{Usuarios:} Redactores (Admin), Estudiantes (Lectores).
        \item \textbf{2. Votaciones del Consejo Estudiantil:} \\ \textit{Problema:} El conteo de votos en papel es lento y poco transparente. \\ \textit{Usuarios:} Comité Electoral (Admin), Estudiantes (Votantes).
        \item \textbf{3. Intercambio de Libros Usados:} \\ \textit{Problema:} Es difícil encontrar quién tiene un libro que ya no usa. \\ \textit{Usuarios:} Estudiantes ofreciendo libros, Estudiantes buscando libros.
        \item \textbf{4. Menú de Cafetería Estudiantil:} \\ \textit{Problema:} Filas largas en el recreo porque los alumnos no saben qué hay. \\ \textit{Usuarios:} Encargado de Cafetería (Admin), Estudiantes (Clientes).
        \item \textbf{5. Torneo Deportivo del Colegio:} \\ \textit{Problema:} Nadie sabe a qué hora juega su equipo ni la tabla de posiciones. \\ \textit{Usuarios:} Organizador (Admin), Jugadores/Espectadores (Usuarios).
    \end{itemize}
    \end{small}
\end{frame}

\begin{frame}{Más Ideas de Proyectos}
    \begin{small}
    \begin{itemize}
        \item \textbf{1. Control de Gastos Personales:} \\ \textit{Problema:} Desconocimiento de en qué se va el dinero (gastos hormiga). \\ \textit{Usuarios:} Cualquier persona (Usuario único).
        \item \textbf{2. Gestor de Tareas y Hábitos:} \\ \textit{Problema:} Desorganización en la rutina diaria y olvido de responsabilidades. \\ \textit{Usuarios:} Estudiantes o Profesionales (Usuario único).
        \item \textbf{3. Catálogo para Emprendedores:} \\ \textit{Problema:} Ventas desorganizadas por WhatsApp y falta de un menú formal. \\ \textit{Usuarios:} Vendedor (Admin), Compradores (Clientes).
        \item \textbf{4. Inventario de Despensa:} \\ \textit{Problema:} Comprar productos repetidos o dejar vencer alimentos por olvido. \\ \textit{Usuarios:} Administrador del hogar (Admin), Miembros de la familia.
        \item \textbf{5. Directorio de Emergencias Locales:} \\ \textit{Problema:} Demora en encontrar números críticos durante una crisis. \\ \textit{Usuarios:} Todos los miembros de una familia o comunidad.
    \end{itemize}
    \end{small}
\end{frame}

\begin{frame}{Próximos Pasos}
    \begin{itemize}
        \item \textbf{Para la próxima clase:}
        \begin{itemize}
            \item Idea de proyecto definida.
            \item Documento inicial con: Problema, Usuarios Objetivo.
            \item Repositorio de GitHub creado con todo el equipo dentro.
        \end{itemize}
    \end{itemize}
    \vspace{1cm}
    \begin{center}
        \Large \textbf{¡Manos a la obra!}
    \end{center}
\end{frame}

\end{document}
